\begin{abstract}
LLM judges now rank models, filter data, and pick the best of $n$ samples, and the standard way to decide whether to trust one is raw agreement with human labels.
That check fails in the worst way: a judge can agree with humans on unoptimized data and still be gamed the moment something optimizes against it.
We introduce a pre-deployment audit: a few dozen blind human labels, graded separately for soundness and for style, reveal whether a judge scores substance or presentation, and yield a verdict of pass, fail, or convict.
We validate the verdicts by running best-of-$n$ selection against every judge and checking who actually gets gamed: the frozen verdicts called the outcome for ten of eleven registered configurations.
The pressure is selection rather than training; the first registered DPO run produced no detectable policy movement (every gain and wedge crossing zero) and was classified post hoc as a manipulation failure, so the predictions went untested, and a stronger rerun carrying a frozen manipulation check is registered and queued.
The audit convicts a deliberately sabotaged frontier judge and a top open reward model, two evaluators that raw agreement cannot tell apart from honest mid-size judges.
The honest pass and the sabotage conviction reproduce significantly on labels from a hypothesis-blind annotator outside the author list; the reward-model conviction reproduces in sign on those labels and is significant on two-rater consensus.
Under pressure, the sabotaged frontier judge inflates its own scoreboard by more than twice the honest judge's inflation while selecting almost as well --- exactly the failure a raw accuracy check cannot see.
The audit costs four to eight annotator-hours of labeling per task and prompt, with one label set serving all thirteen judge configurations here, turning evaluator trust from a post-hoc discovery into a measurement made before deployment.
\end{abstract}
